\section{Tasa útil, latencia y contraste con vectores}
\label{sec:rates-checked}

\paragraph{Coste del camino directo.} Por cada salida de 128 bits, ambos perfiles consumen un bloque
intermedio fresco de 128 bits. Los 256 ceros en el argumento CTR se construyen localmente.
AES/SHA aplica un cifrado AES-256 y comprime un único bloque SHA-256 de 512 bits después del
padding de su mensaje de 384 bits. El perfil CTR ejecuta siete cifrados AES-256 repartidos en dos
claves sucesivas, tres antes de la salida y cuatro para la salida y la actualización posterior.

\begin{center}\small
\begin{tabular}{@{}lrrr@{}}\toprule
Núcleo en software & Mbit/s & latencia p50 ($\mu$s) & latencia p99 ($\mu$s)\\\midrule
AES/SHA & 25,83 & 4,68 & 17,94\\
CTR con inyección en $V$ & 8,91 & 13,69 & 41,18\\\bottomrule
\end{tabular}\end{center}
Cada perfil se midió siete veces con 16\,384 llamadas seriales por repetición. El entorno fue
Linux x86-64 con Python 3.12.14, \texttt{cryptography} 50.0.0 y backend AES OpenSSL 4.0.1. El hash se llama mediante
\texttt{hashlib.sha256}. La tasa es la mediana de las siete tasas de pared. Los percentiles
usan pases separados de latencia individual, después de calentamiento y con entradas sintéticas
precalculadas. El código CTR agrupa las palabras de cada clave en una operación del backend y
forma el estado completo en cada llamada. La medición cubre el intérprete y asignaciones de
objetos. Adquisición, Toeplitz, resiembra, E/S de dispositivo y contención entre consumidores quedan
fuera del tramo cronometrado. La medición aísla el núcleo y omite las pausas entre épocas del
perfil con resiembra cada 1024 peticiones. Los tiempos dependen de plataforma y carga.

\paragraph{Del bloque intermedio a las muestras físicas.} Sea $b$ el número de salidas de
128 bits obtenidas de un lote Toeplitz de $m=128b$ bits, con un selector renovado. El consumidor
ejecuta $b$ llamadas completas al núcleo, una por porción de 128 bits. Si el horizonte $q$ es
múltiplo de $b$ y se destina $2^{-\lambda}$ a la cota conjunta del LHL con selector ideal,
una condición suficiente por lote es
\[
 k_b\ \ge\ m+2\left(\lambda+\log_2(q/b)-1\right).
\]
Resulta de sumar $q/b$ cotas $\tfrac12 2^{(m-k_b)/2}$. Para reservar la mitad del presupuesto
a otras contribuciones se añaden dos bits a $k_b$ y se justifican esas contribuciones por separado.
Con $n=\lceil k_b/h\rceil$ muestras de datos y $d=n+m-1$ bits de selector Toeplitz, supóngase un
bit raw por muestra, $t$ incrementos del selector de $d$ muestras cada uno y $g$ muestras de
guarda. Si cada intento cuesta $n+td+g$ y la probabilidad de aceptación es $p$ en cada intento
independiente, el coste esperado por lote aceptado es $C_b=(n+td+g)/p$. La cota $h$ de entropía
por muestra debe valer en el modelo condicionado a la selección de lotes aceptados, los datos
laterales y el historial. La ley del selector tras $t$ incrementos es otra premisa del modelo.

\begin{center}\small
\begin{tabular}{@{}rrrrrr@{}}\toprule
$b$ & $k_b$ (bits) & $n$ & $d$ & bits por muestra raw & Mbit/s de fuente\\\midrule
1 &318&424&551&0,08134&8,13\\
4 &698&931&1442&0,13176&13,18\\
16&2230&2974&5021&0,15539&15,54\\\bottomrule
\end{tabular}\end{center}
Escenario sintético común a las dos ramas: $q=2^{32}$, $\lambda=64$, $h=0{,}75$, $t=2$,
$g=32$, $p=0{,}99$ y 100 millones de muestras por segundo. La columna de Mbit/s de fuente es la predicción
$10^2(128b/C_b)$, previa a cómputo y consumos auxiliares. La frecuencia de adquisición, la
aceptación y las tasas físicas de las tres filas son exclusivamente valores del modelo.
El experimento de software mide únicamente los dos núcleos criptográficos. En una tubería ideal con etapas
solapadas la menor de las tasas de fuente y cómputo acota la tasa final. Un consumidor serial
suma tiempos de etapas. Arranque, resiembra, calentamiento del selector, recuperación tras
fallos y lotes parcialmente consumidos requieren sumar sus muestras al denominador.
La extensión $b>1$ necesita evidencia de calidad del lote completo y justificar ante AIS la
definición del número intermedio, su partición y los §§338 y 385--390. La fila $b=1$ corresponde
al perfil directo descrito en las dos tuplas. Sus 8,13 Mbit/s corresponden a
424 muestras de datos con presupuesto de selector ideal, antes de asignar semilla.
La sección~\ref{sec:port-rates} usa 427 muestras, reserva error para un selector
real y carga la siembra a cada perfil.

\paragraph{Coste de la siembra y del acondicionamiento.} Para $L$ salidas de 128 bits por época,
contando una aportación inicial o de reseed por época, AES/SHA consume $L+2$ bloques intermedios
en el diseño propuesto. CTR directo consume $L+3$ si un puerto validado entrega directamente
los tres bloques de 128 bits de material de siembra de calidad suficiente. Si este material se
obtiene con las tres evaluaciones CMAC estudiadas, el coste es $L+6$. Usar también CMAC de dos
bloques antes de cada llamada CTR eleva el coste a $2L+6$. Cada etiqueta completa de 128 bits
para el camino externo vetted requiere acreditar 192 bits de entropía de entrada según
§3.2.2.2, además de la hipótesis condicional para el último bloque si se promete ITS.
En RBG3(RS) con AES-256 sin DF, $s=256$ y peticiones de 128 bits, se adquieren por petición
384 bits de material de reseed con full entropy y se generan dos bloques AES de salida antes de
truncar. Esto corresponde a tres bloques de 128 bits desde el puerto full entropy o seis bloques
intermedios en la opción CMAC, más el coste de instanciación. RBG3(XOR) necesita además de su
siembra un bloque físico de full entropy por cada 128 bits entregados. Los bloques físicos empleados
en el XOR de salida y los reservados a la resiembra usan asignaciones disjuntas\cite[§§3.2.2.2, 6.4, 6.5]{nist90c}.

\paragraph{Contraste independiente.} Cuatro trazas oficiales CAVP del NIST para AES-256 sin DF
cubren los cruces con/sin resiembra y con/sin entrada adicional. Verifican $K,V$ tras la
instanciación, tras las resiembras aplicables y tras cada una de dos generaciones, además de los
512 bits devueltos por la segunda llamada. Esas longitudes pertenecen a los vectores oficiales.
Para los perfiles de esta nota se generaron dos llamadas encadenadas de 128 bits por rama con una
segunda transcripción y el ejecutable OpenSSL. Los ensayos incluyen argumento ausente frente a
ceros presentes, padding, acarreo de contador, resiembra AES/SHA con estado anterior, inversa
CMAC y la reconstrucción CTR ante dos estados conocidos. El mutante que omite el \emph{Update}
posterior reproduce la primera salida y falla en la siguiente. Código, trazas, hashes y siete
mediciones por núcleo se incluyen en \texttt{experiments/}.
\href{https://csrc.nist.gov/CSRC/media/Projects/Cryptographic-Algorithm-Validation-Program/documents/drbg/drbgtestvectors.zip}{Archivo CAVP original del NIST}.

\paragraph{Observación de la salida interna (.10 y PTG.3.9).}
Como complemento de las derivaciones computacionales, se publican cuatro trazas
reproducibles con semillas sintéticas públicas. Para cada núcleo se ejecutaron
16 épocas de 1024 peticiones tanto en modo autónomo como con entrada sintética
presente, dos millones de bits por configuración (exactamente \(2^{21}\)).
Se contó la frecuencia de unos para detectar desequilibrio, el XOR de posiciones
adyacentes disjuntas dentro del bloque para detectar dependencia local y el XOR
por posición de dos bloques sucesivos disjuntos para detectar dependencia temporal.
En este orden se observaron los conteos
\[
\begin{array}{lrrr}
 & \text{unos}&\text{XOR intrabloque}&\text{XOR entre bloques}\\
 \text{AES/SHA autónomo} &1047669&523851&525465\\
 \text{AES/SHA con entrada} &1048799&524705&523667\\
 \text{CTR autónomo} &1047842&524124&524240\\
 \text{CTR con entrada} &1048812&524632&523974
\end{array}
\]
La referencia ideal de bits uniformes independientes da medias de \(2^{20}\),
\(2^{19}\) y \(2^{19}\). Para cada suma de \(n\) indicadores, el umbral bilateral
de desviación es \(\sqrt{(n/2)\ln(24/0.01)}\), por Hoeffding y una unión sobre
las doce comparaciones. Las doce quedaron dentro de la banda. Se conserva el
SHA-256 de cada flujo en \path{experiments/diagnostics.json}. El resultado
caracteriza esas ejecuciones de software. La evidencia de la fuente física y
del dispositivo en sus condiciones reales requiere otros datos y controles
\cite[§§157--162, 200 y 209]{ais}.
